%------------------------------------------------------------------------------%
\documentclass[fleqn,utf8,aspectratio=169,12pt]{beamer}

\usepackage{gaal}

\title{Distâncias}

%------------------------------------------------------------------------------%
\begin{document}

\addtitle

%------------------------------------------------------------------------------%
\section{Distância entre Pontos}

%------------------------------------------------------------------------------%
\framedgraphic*{Distância entre Dois Pontos}{F-06-Distancias-pitagoras}{}

%------------------------------------------------------------------------------%
\begin{frame}[t]{Distância entre Dois Pontos}
  Dados
  \[
    A = ( x_1, y_1 )
    \qquad\qquad
    B = ( x_2, y_2 )
  \]
  \pause
  O vetor que liga $A$ a $B$ é
  \[
    \pause
    \overrightarrow{AB} = \vetor{x_2-x_1}{y_2-y_1}
  \]
  \pause
  A distância entre os pontos é a norma do vetor
  \[
    \pause
    d(A, B)
    \pause
    = \norm{\overrightarrow{AB}}
    \pause
    = \sqrt{(x_2-x_1)^{2} + (y_2-y_1)^{2}}
  \]
\end{frame}

%------------------------------------------------------------------------------%
\begin{frame}[t]{Distância no Espaço $\R^3$}
  Dados
  \[
    A = (x_1, y_1, z_1)
    \qquad\qquad
    B = (x_2, y_2, z_2)
  \]
  \pause
  Temos
  \[
    \overrightarrow{AB}
    = \vetortri{x_2-x_1}{y_2-y_1}{z_2-z_1}
  \]
  \pause
  Assim
  \[
    d(A, B)
    \pause
    = \norm{\overrightarrow{AB}}
    \pause
    = \sqrt{ (x_2-x_1)^{2} + (y_2-y_1)^{2} + (z_2-z_1)^{2} }
  \]
\end{frame}

%------------------------------------------------------------------------------%
\section{Propriedades da Distância}

%------------------------------------------------------------------------------%
\begin{frame}{Propriedades}
  Para quaisquer pontos $A$, $B$ e $C$
  \[
    \pause
    d(A, B) \geq 0
  \]

  \[
    \pause
    d(A, B) = 0 \iff A = B
  \]

  \[
    \pause
    d(A, B) = d(B, A)
  \]

  \pause
  Desigualdade triangular
  \[
    \pause
    d(A, C) \leq d(A, B) + d(B, C)
  \]
\end{frame}

%------------------------------------------------------------------------------%
%------------------------------------------------------------------------------%
\begin{question}
  Calcule a distância entre
  \[
    A = (1, 2)
  \]
  \[
    B = (5, 5)
  \]

\end{question}

%------------------------------------------------------------------------------%
\begin{resolution}{Solução}
\begin{columns}[t]
\column{0.5\textwidth}
  \begin{align*}
    \onslide<+->{\overrightarrow{AB} & = B - A           \\}
    \onslide<+->{                    & = (5, 5) - (1, 2) \\}
    \onslide<+->{                    & = (5-1, 5-2)      \\}
    \onslide<+->{                    & = (4, 3)}
  \end{align*}
\column{0.5\textwidth}
  \begin{align*}
    \onslide<+->{d(A, B) & = \norm{\overrightarrow{AB}} \\}
    \onslide<+->{        & = \norm{(4, 3)}              \\}
    \onslide<+->{        & = \sqrt{4^{2}+3^{2}}         \\}
    \onslide<+->{        & = \sqrt{16+9}                \\}
    \onslide<+->{        & = \sqrt{25}                  \\}
    \onslide<+->{        & = 5}
  \end{align*}
\end{columns}
\end{resolution}

%------------------------------------------------------------------------------%
\endinput

%------------------------------------------------------------------------------%
\begin{question}
  Calcule a distância entre
  \[
    A = (-2, 3)
  \]
  \[
    B = (4, -5)
  \]
\end{question}

%------------------------------------------------------------------------------%
\begin{resolution}{Solução}
\begin{columns}[t]
\column{0.5\textwidth}
  \begin{align*}
    \onslide<+->{\overrightarrow{AB} & = B - A             \\}
    \onslide<+->{                    & = (-2, 3) - (4, -5) \\}
    \onslide<+->{                    & = (6, -8)}
  \end{align*}
\column{0.5\textwidth}
  \begin{align*}
    \onslide<+->{d(A, B) & = \norm{\overrightarrow{AB}} \\}
    \onslide<+->{        & = \norm{(6, -8)}             \\}
    \onslide<+->{        & = \sqrt{6^{2}+(-8)^{2}}      \\}
    \onslide<+->{        & = \sqrt{36+64}               \\}
    \onslide<+->{        & = \sqrt{100}                 \\}
    \onslide<+->{        & = 10}
  \end{align*}
\end{columns}
\end{resolution}

%------------------------------------------------------------------------------%
\endinput

%------------------------------------------------------------------------------%
\begin{question}
  Calcule a distância entre
  \[
    A = (1, -2, 3)
  \]
  \[
    B = (4, 2, -1)
  \]
\end{question}

%------------------------------------------------------------------------------%
\begin{resolution}{Solução}
\begin{columns}[t]
\column{0.5\textwidth}
  \begin{align*}
    \onslide<+->{\overrightarrow{AB} & = B - A                   \\}
    \onslide<+->{                    & = (4, 2, -1) - (1, -2, 3) \\}
    \onslide<+->{                    & = (4-1, 2-(-2), -1-3)     \\}
    \onslide<+->{                    & = (3, 4, -4)}
  \end{align*}
\column{0.5\textwidth}
  \begin{align*}
    \onslide<+->{d(A, B) & = \norm{\overrightarrow{AB}}  \\}
    \onslide<+->{        & = \norm{(3, 4, -4)}           \\}
    \onslide<+->{        & = \sqrt{3^{2}+4^{2}+(-4)^{2}} \\}
    \onslide<+->{        & = \sqrt{9+16+16}              \\}
    \onslide<+->{        & = \sqrt{41}}
  \end{align*}
\end{columns}
\end{resolution}

%------------------------------------------------------------------------------%
\endinput

%%------------------------------------------------------------------------------%
\begin{question}
\end{question}

%------------------------------------------------------------------------------%
\begin{resolution}{Solução}
\end{resolution}

%------------------------------------------------------------------------------%
\endinput

% ============================================================
\begin{frame}{Exemplo 4}
Considere
\[
A=(2,-1,4)
\]
e
\[
B=(-1,3,2)
\]

Calcule
\[
\norm{B-A}
\]
e interprete geometricamente o resultado.
\end{frame}

% ============================================================
\begin{frame}{Exemplo 4 --- Solução}
\begin{align*}
B-A
&=
(-1,3,2)-(2,-1,4)
\\
&=
(-3,4,-2)
\end{align*}

Logo,
\begin{align*}
\norm{B-A}
&=
\sqrt{(-3)^{2}+4^{2}+(-2)^{2}}
\\
&=
\sqrt{9+16+4}
\\
&=
\sqrt{29}
\end{align*}

Portanto,
\[
d(A,B)=\sqrt{29}
\]

Geometricamente, $\sqrt{29}$ é o comprimento do segmento que liga $A$ a $B$.
\end{frame}


%------------------------------------------------------------------------------%
%%------------------------------------------------------------------------------%
\begin{question}
\end{question}

%------------------------------------------------------------------------------%
\begin{resolution}{Solução}
\end{resolution}

%------------------------------------------------------------------------------%
\endinput

% ============================================================
\begin{frame}{Exemplo 5}
Considere os pontos
\[
A=(0,0)
\]
\[
B=(3,4)
\]
\[
C=(6,0)
\]

Calcule
\[
d(A,B)
\]
\[
d(B,C)
\]
e
\[
d(A,C)
\]

Verifique a desigualdade triangular.
\end{frame}

% ============================================================
\begin{frame}{Exemplo 5 --- Solução}
Primeiro,
\begin{align*}
\overrightarrow{AB}
&=
(3,4)
\end{align*}

\begin{align*}
d(A,B)
&=
\sqrt{3^{2}+4^{2}}
\\
&=
5
\end{align*}

Agora,
\begin{align*}
\overrightarrow{BC}
&=
(6-3,0-4)
\\
&=
(3,-4)
\end{align*}

\begin{align*}
d(B,C)
&=
\sqrt{3^{2}+(-4)^{2}}
\\
&=
5
\end{align*}
\end{frame}

% ============================================================
\begin{frame}{Exemplo 5 --- Solução}
Finalmente,
\begin{align*}
\overrightarrow{AC}
&=
(6,0)
\end{align*}

\begin{align*}
d(A,C)
&=
\sqrt{6^{2}+0^{2}}
\\
&=
6
\end{align*}

Verificando a desigualdade triangular:
\begin{align*}
d(A,C)
&\leq d(A,B)+d(B,C)
\\
6
&\leq5+5
\\
6
&\leq10
\end{align*}
\end{frame}


%%------------------------------------------------------------------------------%
%\section{Ponto médio}
%
%%------------------------------------------------------------------------------%
%\begin{frame}{Ponto médio de um segmento}
%  Considere
%  \[
%  A=(x_1,y_1)
%  \]
%  e
%  \[
%  B=(x_2,y_2)
%  \]
%
%  O ponto médio $M$ é dado por
%  \[
%  M=
%  \left(
%  \frac{x_1+x_2}{2},
%  \frac{y_1+y_2}{2}
%  \right)
%  \]
%
%  No espaço:
%  \[
%  M=
%  \left(
%  \frac{x_1+x_2}{2},
%  \frac{y_1+y_2}{2},
%  \frac{z_1+z_2}{2}
%  \right)
%  \]
%\end{frame}

%------------------------------------------------------------------------------%
%%------------------------------------------------------------------------------%
\begin{question}
\end{question}

%------------------------------------------------------------------------------%
\begin{resolution}{Solução}
\end{resolution}

%------------------------------------------------------------------------------%
\endinput


% ============================================================
\begin{frame}{Exemplo 6}
Determine o ponto médio do segmento com extremos
\[
A=(-2,4,1)
\]
e
\[
B=(6,-2,5)
\]
\end{frame}

% ============================================================
\begin{frame}{Exemplo 6 --- Solução}
\begin{align*}
M
&=
\left(
\frac{-2+6}{2},
\frac{4+(-2)}{2},
\frac{1+5}{2}
\right)
\\
&=
\left(
\frac{4}{2},
\frac{2}{2},
\frac{6}{2}
\right)
\\
&=
(2,1,3)
\end{align*}
\end{frame}

%%------------------------------------------------------------------------------%
\begin{question}
\end{question}

%------------------------------------------------------------------------------%
\begin{resolution}{Solução}
\end{resolution}

%------------------------------------------------------------------------------%
\endinput

% ============================================================
\begin{frame}{Exemplo 7}
Considere
\[
A=(1,2)
\]
e o ponto médio
\[
M=(4,-1)
\]

Determine o ponto $B$.
\end{frame}

% ============================================================
\begin{frame}{Exemplo 7 --- Solução}
Pela fórmula do ponto médio:
\[
M=
\left(
\frac{x_{A}+x_{B}}{2},
\frac{y_{A}+y_{B}}{2}
\right)
\]

Assim,
\begin{align*}
4
&=
\frac{1+x_{B}}{2}
\\
8
&=
1+x_{B}
\\
x_{B}
&=
7
\end{align*}

Para a segunda coordenada:
\begin{align*}
-1
&=
\frac{2+y_{B}}{2}
\\
-2
&=
2+y_{B}
\\
y_{B}
&=
-4
\end{align*}

Portanto,
\[
B=(7,-4)
\]
\end{frame}

%------------------------------------------------------------------------------%
\begin{question}
  Determine o valor de $k$ para que a distância entre
  \[
    A = (k, 2)
  \]
  \[
    B = (3, -2)
  \]
  seja igual a 5
\end{question}

%------------------------------------------------------------------------------%
\begin{resolution}{Solução}
\begin{columns}[t]
\column{0.5\textwidth}
  \begin{align*}
    \onslide<+->{d(A, B)}
    \onslide<+->{& = \norm{\overrightarrow{AB}}  \\}
    \onslide<+->{& = \sqrt{(3-k)^{2}+(-2-2)^{2}} \\}
    \onslide<+->{& = \sqrt{(3-k)^{2}+16}}
  \end{align*}
\column{0.5\textwidth}
  \begin{align*}
    \onslide<+->{d(A, B)             & = 5     \\}
    \onslide<+->{\sqrt{(3-k)^{2}+16} & = 5     \\}
    \onslide<+->{(3-k)^{2} + 16      & = 25    \\}
    \onslide<+->{(3-k)^{2}           & = 9     \\}
    \onslide<+->{\abs{3-k}           & = 3     \\}
    \onslide<+->{3-k                 & = \pm 3}
  \end{align*}
\end{columns}
\end{resolution}

%------------------------------------------------------------------------------%
\begin{resolution}{Solução}
\begin{columns}[t]
\column{0.3\textwidth}
  \begin{align*}
    \onslide<+->{3-k & = 3 \\}
    \onslide<+->{-k  & = 0 \\}
    \onslide<+->{k   & = 0}
  \end{align*}
\column{0.3\textwidth}
  \begin{align*}
    \onslide<+->{3-k & = -3 \\}
    \onslide<+->{-k  & = -6 \\}
    \onslide<+->{k   & = 6}
  \end{align*}
\end{columns}

  \bigskip
  \onslide<+->{Portanto, \( k = 0 \) ou \( k = 6 \)}
\end{resolution}

%------------------------------------------------------------------------------%
\endinput


%------------------------------------------------------------------------------%
\section{Distância entre Retas Paralelas}

%------------------------------------------------------------------------------%
\framedgraphic*{Distância entre Retas Paralelas}{F-06-Distancias-paralelas-1}{}
\framedgraphic*{Distância entre Retas Paralelas}{F-06-Distancias-paralelas-2}{}
\framedgraphic*{Distância entre Retas Paralelas}{F-06-Distancias-paralelas-3}{}
\framedgraphic*{Distância entre Retas Paralelas}{F-06-Distancias-paralelas-4}{}
\framedgraphic*{Distância entre Retas Paralelas}{F-06-Distancias-paralelas-5}{}
\framedgraphic*{Distância entre Retas Paralelas}{F-06-Distancias-paralelas-6}{}


%------------------------------------------------------------------------------%
\begin{frame}{Distância entre Retas Paralelas}
  A distância entre as \emph{retas paralelas}
  \[
    r\colon \vec{r} = \vec{p} + t\vec{v}
    \qquad\qquad
    s\colon \vec{r} = \vec{q} + \tau\vec{w}
    \qquad\qquad
    \vec{w} = \lambda\vec{v}
  \]
  \pause
  é a norma da componente perpendicular da decomposição de
  \[
    \vec{h} = \vec{q} - \vec{p}
  \]
  \pause
  com relação ao vetor diretor $\vec{v}$
  \[
    \pause
    d(r, s)
    \pause
    = \norm{\vec{h}_\perp}
    \pause
    = \norm{\vec{h} - \projuv{h}{v}}
    \pause
    = \norm{\,\vec{h} - \frac{\;\vec{h}\cdot\vec{v}\;}{\vec{v}\cdot\vec{v}}\vec{v}\,}
  \]
\end{frame}

%------------------------------------------------------------------------------%
\framedgraphic*{Fórmula Alternativa}{F-06-Distancias-paralelogramo-1}{}
\framedgraphic*{Fórmula Alternativa}{F-06-Distancias-paralelogramo-2}{}
\framedgraphic*{Fórmula Alternativa}{F-06-Distancias-paralelogramo-3}{}
\framedgraphic*{Fórmula Alternativa}{F-06-Distancias-paralelogramo-4}{}

%------------------------------------------------------------------------------%
\begin{frame}{Fórmula Alternativa}
  \onslide<+->{Área do paralelogramo definido pelos vetores $\vec{v}$ e $\vec{h}$}
  \[
    \onslide<+->{A = \norm{\vec{h}\times\vec{v}}}
  \]
  \[
    \onslide<+->{A}
    \onslide<+->{= \text{base} \times \text{altura}}
    \onslide<+->{= \norm{\vec{v}} \, d(r, s)}
  \]
  \begin{align*}
    \onslide<+->{\norm{\vec{v}} \, d(r, s) & = \norm{\vec{h}\times\vec{v}} \\}
    \onslide<+->{d(r, s) & =  \frac{\norm{\vec{h}\times\vec{v}}}{\norm{\vec{v}}}}
    \onslide<+->{        \;=\; \frac{\norm{(\vec{q}-\vec{q})\times\vec v}}{\norm{\vec{v}}}}
  \end{align*}
\end{frame}

%------------------------------------------------------------------------------%
%------------------------------------------------------------------------------%
\begin{question}
  Calcule a distância entre as retas
  \[
    r\colon
    \begin{cases}
      x = t   \\
      y = 1+t \\
      z = 2
    \end{cases}
  \]
  \[
    s\colon
    \begin{cases}
      x = 1+\tau \\
      y = 3+\tau \\
      z = 4
    \end{cases}
  \]
\end{question}

%------------------------------------------------------------------------------%
\begin{resolution}{Solução}
  Vetores diretores
  \[
    \vec{w} = \vec{v} = (1, 1, 0)
  \]
  \pause
  Ponto em $r$
  \[
    P = ( 0, 1, 2)
  \]
  \pause
  Ponto em $s$
  \[
    Q = (1, 3, 4)
  \]
\end{resolution}

%------------------------------------------------------------------------------%
\begin{resolution}{Solução}
  \[
    \vec{h}
    \pause
    = Q - P
    \pause
    = (1, 3, 4) - ( 0, 1, 2)
    \pause
    = (1, 2, 2)
  \]
  \[
    \pause
    \vec{h}\times\vec{v}
    \pause
    = \begin{vmatrix}
      \vec{i} & \vec{j} & \vec{k} \\
      1       & 2       & 2       \\
      1       & 1       & 0
    \end{vmatrix}
    \pause
    = (-2, 2, -1)
  \]
  \[
    \pause
    \norm{\vec{h}\times\vec{v}}
    \pause
    = \sqrt{4+4+1}
    \pause
    = 3
  \]
  \[
    \pause
    \norm{\vec{v}}
    \pause
    = \sqrt{1+1}
    \pause
    = \sqrt{2}
  \]
  \[
    \pause
    d(r, s)
    \pause
    = \frac{\norm{(\vec{h})\times\vec v}}{\norm{\vec{v}}}
    \pause
    = \frac{3}{\sqrt{2}}
  \]
\end{resolution}

%------------------------------------------------------------------------------%
\endinput


%------------------------------------------------------------------------------%
\section{Distância de Ponto a Plano}

%------------------------------------------------------------------------------%
\begin{frame}{Distância de um Ponto a um Plano}
  Dados o plano \(\gamma\) e o ponto $P$
  \[
    \gamma\colon ax + by + cz + d = 0
    \qquad\qquad
    P = (x_0, y_0, z_0)
  \]

  \pause
  A distância de $P$ ao plano é
  \[
    \pause
    d(P, \gamma)
    \pause
    = \frac{\, \abs{ax_0 + by_0 + cz_0 + d} \,}{\norm{\vec{n}}}
    \pause
    \qquad\qquad
    \vec{n} = (a, b, c)
  \]
\end{frame}

%------------------------------------------------------------------------------%
%------------------------------------------------------------------------------%
\begin{question}
  Calcule a distância do ponto
  \[
    P = (2, -1, 3)
  \]
  ao plano
  \[
  \gamma\colon 2x - y + 2z - 5 = 0
  \]
\end{question}

%------------------------------------------------------------------------------%
\begin{resolution}{Solução}
  \onslide<+->{Vetor normal do plano}
  \[
    \onslide<+->{\vec{n} = (2, -1, 2)}
  \]
  \[
    \onslide<+->{\norm{\vec{n}}}
    \onslide<+->{= \sqrt{2^{2}+(-1)^{2}+2^{2}}}
    \onslide<+->{= 3}
  \]

  \onslide<+->{Expressão do plano no ponto $P = (2, -1, 3)$}
  \begin{align*}
    \onslide<+->{2x_0 - y_0 + 2z_0 - 5}
    \onslide<+->{& = 2(2) - (-1) + 2(3) - 5 \\}
    \onslide<+->{& = 4 + 1 + 6 - 5          \\}
    \onslide<+->{& = 6}
  \end{align*}
\end{resolution}

%------------------------------------------------------------------------------%
\begin{resolution}{Solução}
  Distância
  \[
    d(P, \gamma)
    \pause
    = \frac{\, \abs{ax_0 + by_0 + cz_0 + d} \,}{\norm{\vec{n}}}
    \pause
    = \frac{\abs{6}}{3}
    \pause
    = 2
  \]
\end{resolution}

%------------------------------------------------------------------------------%
\endinput


%%------------------------------------------------------------------------------%
%\section{Lista Mínima}
%
%%------------------------------------------------------------------------------%
%\begin{frame}{Lista Mínima}
%  \begin{enumerate}
%    \item
%  \end{enumerate}
%
%  \vfill
%  Atenção: A prova é baseada no livro, não nas apresentações
%\end{frame}

%------------------------------------------------------------------------------%
\end{document}
%------------------------------------------------------------------------------%
